\section{Билет 71. Уравнение Ван-дер-Ваальса для реальных газов. Постоянные Ван-дер-Ваальса. Критическая точка. Критические параметры газа (температура, объем и давление). Недостатки уравнения Ван-дер-Ваальса. Изотермы газа Ван-дер-Ваальса.}

\begin{definition}
\textbf{Уравнение Ван-дер-Ваальса} --- одно из простейших уравнений состояния реального газа, учитывающее конечный размер молекул и силы их взаимного притяжения:
\begin{equation}
    \left(p + \frac{a \nu^2}{V^2}\right)(V - \nu b) = \nu R T. \label{eq:vdw_71}
\end{equation}
Для одного моля ($\nu = 1$) уравнение принимает вид: $\left(p + \frac{a}{V^2}\right)(V - b) = R T$.
\end{definition}

\begin{definition}
\textbf{Постоянные Ван-дер-Ваальса} $a$ и $b$ характеризуют конкретное вещество:
\begin{itemize}
    \item $b$ --- \textbf{собственный объём молекул} (исключённый объём). Учитывает, что молекулы имеют конечные размеры и не могут занимать объём меньше $\nu b$.
    \item $a$ --- \textbf{параметр, учитывающий силы межмолекулярного притяжения}. Добавка $\frac{a \nu^2}{V^2}$ называется \textbf{внутренним давлением}; она уменьшает внешнее давление, так как молекулы в поверхностном слое испытывают силу притяжения, направленную внутрь жидкости.
\end{itemize}
\end{definition}

\begin{theorem}[Критическая точка и критические параметры]
\textbf{Критической точкой} называется состояние, в котором исчезает различие между жидкой и газообразной фазами. На изотерме $p(V)$ в критической точке горизонтальная касательная и точка перегиба совпадают:
\begin{equation}
    \left(\frac{\partial p}{\partial V}\right)_T = 0, \quad \left(\frac{\partial^2 p}{\partial V^2}\right)_T = 0. \label{eq:crit_point_71}
\end{equation}
Решая систему \eqref{eq:crit_point_71} совместно с уравнением \eqref{eq:vdw_71}, получаем \textbf{критические параметры}:
\begin{equation}
    V_{\text{кр}} = 3\nu b, \quad p_{\text{кр}} = \frac{a}{27b^2}, \quad T_{\text{кр}} = \frac{8a}{27Rb}. \label{eq:crit_params_71}
\end{equation}
Зная критические параметры, можно экспериментально определить постоянные Ван-дер-Ваальса: $b = \frac{V_{\text{кр}}}{3\nu}$, $a = 3 p_{\text{кр}} V_{\text{кр}}^2$.
\end{theorem}

\begin{remark}
\textbf{Недостатки уравнения Ван-дер-Ваальса:}
\begin{enumerate}
    \item Является приближённым; качественно верно описывает поведение газа, но количественные расчёты (особенно вблизи критической точки и для жидкой фазы) дают значительные погрешности.
    \item Предполагает сферическую симметрию молекул и однородность сил притяжения, что не выполняется для сложных молекул.
    \item Не учитывает ассоциацию молекул (образование димеров, кластеров) и квантовые эффекты при низких температурах.
    \item В области фазового перехода предсказывает физически не реализуемую S-образную кривую (требует замены на горизонтальную площадку по правилу Максвелла).
\end{enumerate}
\end{remark}

\begin{remark}
\textbf{Изотермы газа Ван-дер-Ваальса.}
При $T > T_{\text{кр}}$ изотермы монотонно убывают (похожи на гиперболы). При $T < T_{\text{кр}}$ изотермы имеют максимум и минимум, образуя S-образную кривую. Участки между максимумом и минимумом соответствуют метастабильным состояниям (перегретая жидкость, пересыщенный пар), а участок с положительной производной $\partial p/\partial V > 0$ физически не реализуем (термодинамически неустойчивое состояние).
\end{remark}
